% GENERATED FILE. Edit canonical lessons, then run npm run generate:books.
% canonical-source-hash: fnv1a64:c020a78bea7e731c
% canonical-lessons: TE-C265-pellikuturu, TE-C265-janam, TE-C265-devudu, TE-C265-donga, TE-C265-jantuvu

\chapter{Bride, People, God, Thief, Animal}
\label{ch:te-bride-people-god-thief-animal}
\hlchaptermodality{265}

\begin{chapteropening}
\textbf{By the end of this chapter:} \emph{I can say a bride, people, God, a thief and an animal.}

Complete the last lesson of chapter 265: I can say a bride, people, God, a thief and an animal.
\end{chapteropening}

\section[\texorpdfstring{\te{పెళ్ళికూతురు} (peḷḷikūturu)}{పెళ్ళికూతురు (peḷḷikūturu)}]{\te{పెళ్ళికూతురు} (peḷḷikūturu) — a bride}

\label{lesson:TE-C265-pellikuturu}



\begin{quote}
Before the new one: say the Telugu for a newspaper, then the Telugu for a number.
\end{quote}

\subsection*{You'll want to know: \te{పెళ్ళికూతురు}}
\textbf{\te{పెళ్ళికూతురు}} — \emph{peḷḷikūturu} — \textquotedblleft{}a bride\textquotedblright{}.

\begin{grammarlens}[title={What you've built}]
One more word for this chapter.
\end{grammarlens}

\subsection*{Your turn}
\begin{itemize}
\raggedright
  \item \emph{Say it:} \emph{peḷḷikūturu}
  \item \emph{Say it:} \emph{peḷḷikūturu}, once more
  \item \emph{Say it:} say \emph{peḷḷikūturu}
  \item \emph{Recall:} say the Telugu for a newspaper, then the Telugu for a number, then say \emph{peḷḷikūturu} again
\end{itemize}

\begin{tcolorbox}[breakable,colback=teal!4,colframe=teal!35!black,title={Before you move on}]
What does \te{పెళ్ళికూతురు} mean? (\textquotedblleft{}A bride\textquotedblright{}.) Say it once more.
\end{tcolorbox}
\section[\texorpdfstring{\te{జనం} (janaṁ)}{జనం (janaṁ)}]{\te{జనం} (janaṁ) — people, a crowd}

\label{lesson:TE-C265-janam}



\begin{quote}
Before the new one: say the Telugu for a number, then the Telugu for a bride.
\end{quote}

\subsection*{You'll want to know: \te{జనం}}
\textbf{\te{జనం}} — \emph{janaṁ} — \textquotedblleft{}people, a crowd\textquotedblright{}.

\begin{grammarlens}[title={What you've built}]
One more word for this chapter.
\end{grammarlens}

\subsection*{Your turn}
\begin{itemize}
\raggedright
  \item \emph{Say it:} \emph{janaṁ}
  \item \emph{Say it:} \emph{janaṁ}, once more
  \item \emph{Say it:} say \emph{janaṁ}
  \item \emph{Recall:} say the Telugu for a number, then the Telugu for a bride, then say \emph{janaṁ} again
\end{itemize}

\begin{tcolorbox}[breakable,colback=teal!4,colframe=teal!35!black,title={Before you move on}]
What does \te{జనం} mean? (\textquotedblleft{}People, a crowd\textquotedblright{}.) Say it once more.
\end{tcolorbox}
\section[\texorpdfstring{\te{దేవుడు} (dēvuḍu)}{దేవుడు (dēvuḍu)}]{\te{దేవుడు} (dēvuḍu) — God, a god}

\label{lesson:TE-C265-devudu}



\begin{quote}
Before the new one: say the Telugu for a bride, then the Telugu for people.
\end{quote}

\subsection*{You'll want to know: \te{దేవుడు}}
\textbf{\te{దేవుడు}} — \emph{dēvuḍu} — \textquotedblleft{}God, a god\textquotedblright{}.

\begin{grammarlens}[title={What you've built}]
One more word for this chapter.
\end{grammarlens}

\subsection*{Your turn}
\begin{itemize}
\raggedright
  \item \emph{Say it:} \emph{dēvuḍu}
  \item \emph{Say it:} \emph{dēvuḍu}, once more
  \item \emph{Say it:} say \emph{dēvuḍu}
  \item \emph{Recall:} say the Telugu for a bride, then the Telugu for people, then say \emph{dēvuḍu} again
\end{itemize}

\begin{tcolorbox}[breakable,colback=teal!4,colframe=teal!35!black,title={Before you move on}]
What does \te{దేవుడు} mean? (\textquotedblleft{}God, a god\textquotedblright{}.) Say it once more.
\end{tcolorbox}
\section[\texorpdfstring{\te{దొంగ} (doṅga)}{దొంగ (doṅga)}]{\te{దొంగ} (doṅga) — a thief}

\label{lesson:TE-C265-donga}



\begin{quote}
Before the new one: say the Telugu for people, then the Telugu for God.
\end{quote}

\subsection*{You'll want to know: \te{దొంగ}}
\textbf{\te{దొంగ}} — \emph{doṅga} — \textquotedblleft{}a thief\textquotedblright{}.

\begin{grammarlens}[title={What you've built}]
One more word for this chapter.
\end{grammarlens}

\subsection*{Your turn}
\begin{itemize}
\raggedright
  \item \emph{Say it:} \emph{doṅga}
  \item \emph{Say it:} \emph{doṅga}, once more
  \item \emph{Say it:} say \emph{doṅga}
  \item \emph{Recall:} say the Telugu for people, then the Telugu for God, then say \emph{doṅga} again
\end{itemize}

\begin{tcolorbox}[breakable,colback=teal!4,colframe=teal!35!black,title={Before you move on}]
What does \te{దొంగ} mean? (\textquotedblleft{}A thief\textquotedblright{}.) Say it once more.
\end{tcolorbox}
\section[\texorpdfstring{\te{జంతువు} (jantuvu)}{జంతువు (jantuvu)}]{\te{జంతువు} (jantuvu) — an animal}

\label{lesson:TE-C265-jantuvu}



\begin{quote}
Before the new one: say the Telugu for God, then the Telugu for a thief.
\end{quote}

\subsection*{You'll want to know: \te{జంతువు}}
\textbf{\te{జంతువు}} — \emph{jantuvu} — \textquotedblleft{}an animal\textquotedblright{}.

\begin{grammarlens}[title={What you've built}]
One more word for this chapter.
\end{grammarlens}

\subsection*{Your turn}
\begin{itemize}
\raggedright
  \item \emph{Say it:} \emph{jantuvu}
  \item \emph{Say it:} \emph{jantuvu}, once more
  \item \emph{Say it:} say \emph{jantuvu}
  \item \emph{Recall:} say the Telugu for God, then the Telugu for a thief, then say \emph{jantuvu} again
\end{itemize}

\begin{tcolorbox}[breakable,colback=teal!4,colframe=teal!35!black,title={Before you move on}]
What does \te{జంతువు} mean? (\textquotedblleft{}An animal\textquotedblright{}.) Say it once more.
\end{tcolorbox}
